
En
posant @@M0@@, on se ramène à calculer le résidu de :
@@M1@@
c'est-à-dire le coefficient de @@M2@@ dans @@M3@@. Pour @@M4@@, il est nul, et pour @@M5@@, c'est @@M6@@.

Ainsi s'achève la démonstration du théorème de Riemann--Roch.

\subsection*{BIBLIOGRAPHIE}
\begin{enumerate}[label={[\arabic*]}]
\item A. Borel et J.P. Serre, Le théorème de Riemann--Roch, Bull. Soc. Math. de France, t. 86, 1958.
\item A. Grothendieck, Techniques de construction et théorèmes d'existence en Géométrie Algébrique, Sém. Bourbaki, 1960-61.
\end{enumerate}




